\begin{question}[ID=1.1, subtitle={Teorema di Rolle per $f(x)=x^2-4x+3$ su $[1,3]$}, tags={rolle, teoremi}, difficulty=facile]
Verificare se la funzione \(f(x) = x^2 - 4x + 3\) soddisfa le ipotesi del teorema di Rolle nell'intervallo \([1;\ 3]\) e, in caso affermativo, trovare il punto \(c\) in cui \(f'(c) = 0\).
\begin{solution}
$f(1)=0=f(3)$, $f$ continua e derivabile su $[1,3]$. $f'(x)=2x-4=0 \Rightarrow c=2$
\end{solution}
\end{question}

%-----------------------------%

\begin{question}[ID=1.2, subtitle={Teorema di Rolle per $\sin(x)$ su $[0,\pi]$}, tags={rolle, teoremi}, difficulty=facile]
Sia \(f(x) = \sin(x)\). Verificare che \(f(0) = f(\pi)\) e trovare un valore \(c \in (0;\ \pi)\) tale che \(f'(c) = 0\).
\begin{solution}
$f(0)=0=f(\pi)$. $f'(x)=\cos(x)=0 \Rightarrow c=\dfrac{\pi}{2}$
\end{solution}
\end{question}

%-----------------------------%

\begin{question}[ID=1.3, subtitle={Teorema di Rolle per $f(x)=x^3-x$ su $[-1,1]$}, tags={rolle, teoremi}, difficulty=media]
Considera la funzione \(f(x) = x^3 - x\) nell'intervallo \([-1;\ 1]\). Mostrare che soddisfa le ipotesi del teorema di Rolle e determinare il valore di \(c\) in \((-1, 1)\) per cui \(f'(c) = 0\).
\begin{solution}
$f(-1)=0=f(1)$. $f'(x)=3x^2-1=0 \Rightarrow c=\pm\dfrac{1}{\sqrt{3}}$
\end{solution}
\end{question}

%-----------------------------%

\begin{question}[ID=1.4, subtitle={Teorema di Rolle per $f(x)=e^{-x}\sin(x)$ su $[0,\pi]$}, tags={rolle, teoremi}, difficulty=media]
Applica il teorema di Rolle alla funzione \(f(x) = e^{-x}\sin(x)\) nell'intervallo \([0;\ \pi]\). Trovare il punto \(c\) in cui la derivata si annulla.
\begin{solution}
$f(0)=0=f(\pi)$. $f'(x)=e^{-x}(\cos x - \sin x)=0 \Rightarrow \tan c = 1 \Rightarrow c=\dfrac{\pi}{4}$
\end{solution}
\end{question}

%-----------------------------%

\begin{question}[ID=1.5, subtitle={Teorema di Rolle per $f(x)=x^2-5x$ su $[1,4]$}, tags={rolle, teoremi}, difficulty=facile]
Applica il teorema di Rolle alla funzione \(f(x) = x^2 - 5x\), nell'intervallo \([1;\ 4]\). Trovare il punto \(c\) in cui la derivata si annulla.
\begin{solution}
$f(1)=-4=f(4)$. $f'(x)=2x-5=0 \Rightarrow c=\dfrac{5}{2}$
\end{solution}
\end{question}

%-----------------------------%

\begin{question}[ID=1.6, subtitle={Teorema di Rolle per $f(x)=x^2-5x$ su $[1,4]$ — variante}, tags={rolle, teoremi}, difficulty=facile]
Applica il teorema di Rolle alla funzione \(f(x) = x^2 - 5x\), nell'intervallo \([1;\ 4]\). Trovare il punto \(c\) in cui la derivata si annulla.
\begin{solution}
$f(1)=-4=f(4)$. $f'(x)=2x-5=0 \Rightarrow c=\dfrac{5}{2}$
\end{solution}
\end{question}

%-----------------------------%

\begin{question}[ID=1.7, subtitle={Teorema di Rolle non applicabile: $f(x)=1/(x-1)$ su $[1/2,3]$}, tags={rolle, teoremi, non applicabile}, difficulty=media]
Applica il teorema di Rolle alla funzione \(f(x) = \dfrac{1}{x-1}\), nell'intervallo \(\bigl[\tfrac{1}{2};\ 3\bigr]\). Trovare il punto \(c\) in cui la derivata si annulla.
\begin{solution}
Teorema non applicabile: $f$ è discontinua in $x=1 \in \bigl[\tfrac{1}{2}, 3\bigr]$.
\end{solution}
\end{question}

%-----------------------------%

\begin{question}[ID=1.8, subtitle={Teorema di Rolle non applicabile: $f(x)=4x^2-2x$ su $[-1,3]$}, tags={rolle, teoremi, non applicabile}, difficulty=facile]
Applica il teorema di Rolle alla funzione \(f(x) = 4x^2 -2x\), nell'intervallo \([-1;\ 3]\). Trovare il punto \(c\) in cui la derivata si annulla.
\begin{solution}
Teorema non applicabile: $f(-1)=6 \ne f(3)=30$.
\end{solution}
\end{question}

%-----------------------------%

\begin{question}[ID=2.1, subtitle={Parametro $a$ per Rolle: $ax^3+(a-1)x$ su $[-1,2]$}, tags={rolle, teoremi, parametrico}, difficulty=difficile]
Trova i valori dei parametri in modo che la seguente funzione $f(x) = ax^3 + (a-1)x$ verifichi il teorema di Rolle nell'intervallo $[-1;\ 2]$. Determina il punto (o i punti) la cui esistenza è assicurata dal teorema.
\begin{solution}
$a=\dfrac{1}{4},\ c=1$
\end{solution}
\end{question}

%-----------------------------%

\begin{question}[ID=2.2, subtitle={Parametro $a$ per Rolle: $(a+2)x^4-2ax+4$ su $[-1,1]$}, tags={rolle, teoremi, parametrico}, difficulty=difficile]
Trova i valori dei parametri in modo che la seguente funzione $f(x) = (a+2)x^4 -2ax +4$ verifichi il teorema di Rolle nell'intervallo $[-1;\ 1]$. Determina il punto (o i punti) la cui esistenza è assicurata dal teorema.
\begin{solution}
$a=0,\ c=0$
\end{solution}
\end{question}
